\devoirDS{6}{Bilan annuel et chapitres 19-20}{Un circuit série/6,Une acquisition de température/7,Une solution au laboratoire/7}
\begin{blocLivre}{ds-6-1}\begin{exoBac}{6}{Un circuit série}
Un générateur de 9,0 V alimente deux résistances en série de 100 et 200 ohms.
\schema{serie}
\begin{questions}
\item Calculer l’intensité grâce aux lois d’Ohm et des mailles.
\item Calculer les tensions aux bornes des deux résistances.
\item Indiquer comment mesurer l’intensité et la tension de $R_2$ sans créer de court-circuit.
\end{questions}
\end{exoBac}\end{blocLivre}
\begin{blocLivre}{ds-6-2}\begin{exoBac}{7}{Une acquisition de température}
Le pont utilise $U_0=5,00$ V et $R_0=10,0$ kΩ. Le convertisseur fictif renvoie 0 pour 0 V et 1000 pour 5,00 V. Le code mesuré vaut 500. On donne $R_c=R_0U/(U_0-U)$ et $R_c=20,0-0,400\theta$ en kΩ et °C, de 0 à 40 °C.
\schema{pont-capteur}
\begin{questions}
\item Calculer la tension U.
\item Calculer la résistance puis la température.
\item Une alerte à 30 °C doit-elle se déclencher ?
\item Citer deux vérifications avant d’utiliser cette chaîne sur un autre capteur.
\end{questions}
\end{exoBac}\end{blocLivre}
\begin{blocLivre}{ds-6-3}\begin{exoBac}{7}{Une solution au laboratoire}
On veut 100,0 mL à 4,00 g/L à partir d’une mère à 40,0 g/L. Un chauffage ultérieur au laboratoire laisse une seule tache de produit S au même niveau que la référence N sur une CCM menée dans les mêmes conditions.
\schema{dilution}
\begin{questions}
\item Calculer le prélèvement et choisir la verrerie.
\item Expliquer le reflux qui peut être utilisé pour une synthèse.
\item La tache commune suffit-elle à garantir une pureté absolue ?
\item Une température augmente pendant la transformation dans un montage suffisamment isolé : interpréter et citer une limite de l’observation.
\end{questions}
\end{exoBac}\end{blocLivre}

